\documentclass[11pt]{article}

\usepackage[margin=2.8cm]{geometry}
\usepackage{amsmath,amssymb,amsthm}
\usepackage{stmaryrd}
\usepackage{enumitem}
\usepackage[hidelinks]{hyperref}

\title{A Formalization of ``An Expressive Trace Logic for Recursive
  Programs''\\[0.3em]
  \large How the formalization differs from the paper}
\author{}
\date{}

% ---------------------------------------------------------------------------
% Notation

\newcommand{\sem}[1]{[\![ #1 ]\!]}
\newcommand{\fsem}[1]{|\!|#1|\!|}
\newcommand{\Str}{{\mathcal{S}_{\mathit{tr}}}}
\newcommand{\Ssos}{{\mathcal{S}_{\mathit{sos}}}}
\newcommand{\chop}{\smallfrown}
\newcommand{\Id}{\mathsf{Id}}
\newcommand{\Sb}[2]{\mathsf{Sb}^{#1}_{#2}}
\newcommand{\dup}{\sharp}
\newcommand{\dom}{\mathrm{dom}}
\newcommand{\stf}{\mathrm{stf}}
\newcommand{\State}{\mathit{State}}
\newcommand{\kw}[1]{\mathbf{#1}}
\newcommand{\skipS}{\kw{skip}}
\newcommand{\call}[1]{#1()}
\newcommand{\Tr}{\State^{*}}
\newcommand{\pow}[1]{2^{#1}}
\newcommand{\lfp}{\mathrm{lfp}}
% A numbered item of the paper [GH], e.g. \GH{Definition~5.4}.
\newcommand{\GH}[1]{#1 of GH}

\theoremstyle{plain}
\newtheorem{theorem}{Theorem}[section]
\newtheorem{lemma}[theorem]{Lemma}
\newtheorem{proposition}[theorem]{Proposition}
\newtheorem{corollary}[theorem]{Corollary}
\theoremstyle{definition}
\newtheorem{definition}[theorem]{Definition}
\newtheorem{example}[theorem]{Example}
\theoremstyle{remark}
\newtheorem{remark}[theorem]{Remark}

% Rules of the calculus and of the SOS, generated by Ott from the same
% source as the formalization. The rule label is the optional argument.
\input{ott-rules.tex}
\renewcommand{\ottdrule}[4][]{{\displaystyle
    \frac{\begin{array}{c}#2\end{array}}{#3}\ \textsc{(#1)}}}
\renewcommand{\ottpremise}[1]{#1 \\}

\begin{document}
\maketitle

\section{Introduction}

This note accompanies a HOL4 formalization of parts of the paper
\emph{An Expressive Trace Logic for Recursive Programs} by Gurov and
H\"ahnle~\cite{GH} (henceforth GH). The formalization covers the syntax and
semantics of the programming language Rec, the trace logic, strongest trace
formulas, and the proof calculus. Its main results are soundness of the
calculus (\GH{Theorem~5.6}), the characterisation of strongest trace formulas
(Theorem~4.9), the correctness of the trace semantics with respect to the
SOS (Theorem~3.2), and the example derivation of Figure~6.

The formalization follows the paper closely. This note describes where it
differs, written as if the formalization were a paper: which definitions are
made more precise, which statements change, and how some proofs are done
differently. Results that the formalization proves as stated in the paper,
under the conventions of Section~\ref{sec:prelim}, are mentioned only
briefly. The noteworthy changes and additions are:
\begin{itemize}
\item \emph{Validity of sequents.} With \GH{Definition~5.4}, the calculus
  is not sound for non-empty antecedents
  (Proposition~\ref{prop:counter}). The corrected definition interprets the
  procedures in the domain of the antecedent by their hypotheses
  (Definition~\ref{def:valid}). It agrees with the paper's for empty
  antecedents, and soundness holds for all sequents
  (Section~\ref{sec:validity}).
\item \emph{Semantics relative to a set of procedures}
  (Definition~\ref{def:str}). It makes precise what the (Call) rule assumes,
  and is used both for validity and for strongest trace formulas.
\item \emph{Structural rules and side conditions of the calculus}
  (Section~\ref{sec:calculus}). The paper omits the structural rules; the
  formalization gives them explicitly. Some must preserve the domain of the
  antecedent, and a left rule for $\lor$ is unsound. (Unfold) requires a
  closed fixed-point formula, which makes plain substitution sufficient
  (Lemma~\ref{lem:subst}).
\item \emph{Proof structure.} The (Call) case of soundness is a
  fixed-point induction over all procedures at once, instead of an argument
  for one procedure plus Beki\v{c}'s principle (Lemma~\ref{lem:call}). The
  proof of \GH{Theorem~4.9} does not use modal equation systems; it goes
  through a generalization to open formulas (Lemma~\ref{lem:stf-open}) and
  a single step of Beki\v{c}'s principle (Lemma~\ref{lem:bekic}).
\item \emph{Figure~6 in full} (Section~\ref{sec:example}). The two (Cons)
  steps that the paper justifies informally follow from lemmas about
  unfolding strongest trace formulas (Lemma~\ref{lem:unfold}). A shorter
  derivation needs neither of them (Remark~\ref{rem:short}).
\end{itemize}

Each numbered definition and result in this note corresponds to a HOL4
definition or theorem. The file \texttt{README.md} of the formalization
lists them by the numbers used here. References such as
``\GH{Definition~5.4}'' are to the paper.

\section{Programs and their semantics}\label{sec:prelim}

\paragraph{Conventions.} States are total functions from variable names to
integers. Arithmetic and Boolean expressions are fixed concrete grammars;
Boolean expressions have $=$, $\leq$, $\lnot$ and $\land$, so for example
$x > 0$ is written $\lnot(x \leq 0)$. A trace is a finite sequence of
states, and $\Tr$ denotes the set of traces. The paper uses non-empty
sequences $\State^{+}$; nothing below depends on excluding the empty
sequence.

\begin{definition}[Procedure tables]
  A procedure table $T$ is a finite list of declarations $m\,\{S_m\}$. Lookup
  returns the first declaration of $m$, so duplicate names need no
  well-formedness condition. We write $m\,\{S_m\} \in T$ when lookup of $m$
  returns $S_m$, and $M$ for the set of declared names. $T$ is
  \emph{variable-free} if no $S_m$ contains a statement variable.
\end{definition}

Statements may contain statement variables $Y$ (\GH{Section~5}). A procedure
$m$ and the statement variable $Y_m$ associated with it (\GH{Section~5.1.1})
are identified by their name: $Y_m$ \emph{is} $m$. An interpretation of
statement variables is a function $I$ from names to $\pow{\Tr}$; an
interpretation of procedures $\rho$ has the same type.

\begin{definition}[Trace semantics, relative to a set of procedures]
  \label{def:str}
  The helper semantics $\Str'\sem{S}_{\rho,I}$ is defined by the equations of
  \GH{Figure~2}, with $\Str'\sem{\call{m}}_{\rho,I} = \rho(m)$ and
  $\Str'\sem{Y}_{\rho,I} = I(Y)$. Let $D$ be a set of procedure names, and
  define the transformer $H^D_I$ on interpretations of procedures by
  \[
    H^D_I(\rho)(m) =
    \begin{cases}
      \dup I(Y_m) & \text{if } m \in D,\\
      \dup \Str'\sem{S_m}_{\rho,I} & \text{if } m \notin D \text{ and } m\,\{S_m\} \in T,\\
      \emptyset & \text{otherwise.}
    \end{cases}
  \]
  Let $\rho^D_I$ be its least fixed point, and define
  $\Str^D_I\sem{S} = \Str'\sem{S}_{\rho^D_I,I}$. We write
  $\Str_I\sem{S}$ for $\Str^{\emptyset}_I\sem{S}$; this is the semantics of
  \GH{Definition~3.1}, relativized on $I$ as described in \GH{Section~5}.
\end{definition}

The set $D$ is new. In $\Str^D_I$, a procedure $m \in D$ behaves as
$\skipS; Y_m$, that is, as the hypothesis about $m$ recorded in an
antecedent. This is what the (Call) rule assumes, and it is needed for
validity of sequents (Section~\ref{sec:validity}) and for the proof of
Theorem~\ref{thm:stf}. Undeclared procedures denote $\emptyset$; the paper
assumes that every called procedure is declared.

Least fixed points, here and for $\mu$-formulas, are defined as the
intersection of all pre-fixed points (Knaster--Tarski). Only monotonicity is
needed; the continuity argument of \GH{Section~3} is not used.

\begin{theorem}[\GH{Theorem~3.2}]\label{thm:sos}
  If $T$ and $S$ are variable-free, then $\Str_I\sem{S} = \Ssos\sem{S}$ for
  every $I$.
\end{theorem}

The SOS and the induced traces are defined as in \GH{Figure~1 and
Definition~2.6}; Figure~\ref{fig:sos} in the appendix shows them as
formalized. The SOS has no rule for statement variables, hence the
restriction to variable-free programs. The inclusion
$\Ssos\sem{S} \subseteq \Str_I\sem{S}$ is by induction on the length of the
execution. For the converse, the traces induced by the SOS for the calls
$\call{m}$ form a pre-fixed point of $H^{\emptyset}_I$, so they contain
$\rho^{\emptyset}_I$ by leastness.

\section{Trace formulas}

Formulas are as in \GH{Definition~4.1}. State formulas $p$ are Boolean
expressions. Besides $\Id$ and $\Sb{a}{x}$, relation symbols are interpreted
by a fixed parameter. The semantics $\fsem{\phi}_V$ is that of \GH{Figure~3}, with
$\fsem{\mu X.\phi}_V$ the least pre-fixed point. For an
arbitrary formula we write $\fsem{\phi}$ for $\fsem{\phi}_{V_\emptyset}$,
where $V_\emptyset(X) = \emptyset$ for all $X$; for closed formulas this
agrees with the paper. Substitution $\phi[\psi/X]$ is the plain syntactic
substitution, which stops at binders for $X$ but does not rename bound
variables.

\begin{lemma}[Substitution]\label{lem:subst}
  If $\psi$ is closed, then
  $\fsem{\phi[\psi/X]}_V = \fsem{\phi}_{V[X \mapsto \fsem{\psi}]}$.
  Consequently, if $\mu X.\phi$ is closed, then
  $\fsem{\phi[\mu X.\phi/X]} = \fsem{\mu X.\phi}$ (\GH{Proposition~4.3}).
\end{lemma}

For a closed $\psi$, nothing can be captured, so plain substitution agrees
with capture-avoiding substitution up to renaming of bound variables. The
calculus only substitutes closed formulas (rule (Unfold), with an explicit
side condition, see Section~\ref{sec:calculus}). Capture-avoiding
substitution would only be needed for rules that substitute open formulas.

\section{Strongest trace formulas}

\begin{definition}[\GH{Definition~4.6}]\label{def:stf}
  $\stf(X, S)$ is defined by the equations of \GH{Figure~4}, with
  $\stf(X, Y) = Y$ for a statement variable $Y$ (read as a recursion
  variable). As for statement variables, the recursion variable $X_m$ is
  identified with the name $m$. The recursion terminates: in the case of a
  call, either the number of declared procedures not in $X$ decreases, or it
  stays the same and the statement gets smaller. We write $\stf(S)$ for
  $\stf(\emptyset, S)$.
\end{definition}

Because $X_m$, $Y_m$ and $m$ are identified, a valuation $V$ of recursion
variables is also an interpretation of statement variables. This makes the
following generalization of \GH{Theorem~4.9} to open formulas possible. It
replaces the argument via modal equation systems in \GH{Appendix~A.2}.

\begin{lemma}[One step of Beki\v{c}'s principle]\label{lem:bekic}
  Let $T$ be variable-free, $m \notin D$ and $m\,\{S_m\} \in T$. Then
  \[
    \lfp\big(\gamma \mapsto \Str^{D \cup \{m\}}_{I[Y_m \mapsto \gamma]}\sem{S_m}\big)
    = \Str^D_I\sem{S_m}.
  \]
\end{lemma}

\begin{proof}[Proof sketch]
  Write $G = \Str^D_I\sem{S_m}$. For $\subseteq$, it suffices that $G$ is a
  pre-fixed point. With $I' = I[Y_m \mapsto G]$, the interpretation
  $\rho^D_I$ is a pre-fixed point of $H^{D \cup \{m\}}_{I'}$, since
  $\rho^D_I(m) = \dup G = \dup I'(Y_m)$. So
  $\rho^{D \cup \{m\}}_{I'} \sqsubseteq \rho^D_I$, and
  $\Str^{D \cup \{m\}}_{I'}\sem{S_m} \subseteq G$ by monotonicity. For
  $\supseteq$, let $\gamma$ be any pre-fixed point and
  $I' = I[Y_m \mapsto \gamma]$. Then $\rho^{D \cup \{m\}}_{I'}$ is a
  pre-fixed point of $H^D_I$, so it contains $\rho^D_I$, and hence
  $G \subseteq \gamma$.
\end{proof}

\begin{lemma}[Open strongest trace formulas]\label{lem:stf-open}
  Let $T$ be variable-free, and let $V(m) = \emptyset$ for all undeclared
  $m \notin X$. Then $\fsem{\stf(X, S)}_V = \Str^X_V\sem{S}$.
\end{lemma}

\begin{proof}[Proof sketch]
  By induction along the recursion of Definition~\ref{def:stf}. For a call
  $\call{m}$ with $m \in X$, both sides are $\dup V(m)$. For a call with
  $m \notin X$ and $m\,\{S_m\} \in T$, the left-hand side is
  $\dup\,\lfp(\gamma \mapsto \fsem{\stf(X \cup \{m\}, S_m)}_{V[m \mapsto
    \gamma]})$; by the induction hypothesis and Lemma~\ref{lem:bekic} this is
  $\dup\Str^X_V\sem{S_m} = \rho^X_V(m)$.
\end{proof}

\begin{theorem}[\GH{Theorem~4.9}]\label{thm:stf}
  If $T$ and $S$ are variable-free, then
  $\fsem{\stf(S)} = \Str_I\sem{S} = \Ssos\sem{S}$ for every $I$.
\end{theorem}

The next two lemmas are used to justify the (Cons) steps of Figure~6
(Section~\ref{sec:example}). The first says that the relative semantics
agrees with the ordinary one when the hypotheses are exact.

\begin{lemma}[Exact hypotheses]\label{lem:exact}
  Let $T$ be variable-free. If for every $m \in D$ we have
  $m\,\{S_m\} \in T$ and $I(Y_m) = \Str_I\sem{S_m}$, then
  $\Str^D_I\sem{S} = \Str_I\sem{S}$.
\end{lemma}

\begin{lemma}[Unfolding strongest trace formulas]\label{lem:unfold}
  Let $T$ be variable-free, $n\,\{S_n\} \in T$, and let
  $\phi_n = \mu X_n.\stf(\{n\}, S_n)$ be closed. Then:
  \begin{enumerate}[label=(\alph*)]
  \item $\fsem{\stf(\{n\}, S)[\phi_n/X_n]} = \Str_I\sem{S}$ for every
    variable-free $S$;
  \item $\fsem{(\mu X_m.\stf(\{m, n\}, S_m))[\phi_n/X_n]} =
    \fsem{\mu X_m.\stf(\{m\}, S_m)}$ for every $m \neq n$ with
    $m\,\{S_m\} \in T$.
  \end{enumerate}
\end{lemma}

Both parts follow from Lemmas~\ref{lem:subst}, \ref{lem:bekic},
\ref{lem:stf-open} and~\ref{lem:exact}, since $\fsem{\phi_n} =
\Str_I\sem{S_n}$.

\begin{example}[\GH{Examples~4.7 and~4.8}]\label{ex:stf}
  Evaluating Definition~\ref{def:stf} yields the formulas of the paper,
  with recursion variables named after the procedures. Chop is nested to the
  right, as the definition builds it: $\Id\chop\Sb{x-1}{x}\chop\Id\chop X$
  stands for $\Id\chop(\Sb{x-1}{x}\chop(\Id\chop X))$. In Example~4.8,
  $x > 0$ is $\lnot(x \leq 0)$, and so $x \leq 0$ appears as
  $\lnot\lnot(x \leq 0)$.
\end{example}

\section{The proof calculus}\label{sec:calculus}

\begin{figure}[t]
  \small
  \[
    \ottdrulecalcXXskip{Skip} \qquad
    \ottdrulecalcXXassign{Assign} \qquad
    \ottdrulecalcXXseq{Seq}
  \]
  \[
    \ottdrulecalcXXif{If} \qquad
    \ottdrulecalcXXunfold{Unfold}
  \]
  \[
    \ottdrulecalcXXcall{Call}
  \]
  \[
    \ottdrulecalcXXcons{Cons} \qquad
    \ottdrulecalcXXconsL{ConsL}
  \]
  \[
    \ottdrulecalcXXax{Ax} \qquad
    \ottdrulecalcXXweak{Weak} \qquad
    \ottdrulecalcXXcut{Cut}
  \]
  \[
    \ottdrulecalcXXandR{$\land$R} \qquad
    \ottdrulecalcXXandL{$\land$L}
  \]
  \caption{The proof calculus as formalized (generated from the Ott
    source). The procedure table $T$ and the entailment oracle $\models$
    are parameters of the calculus.}
  \label{fig:calculus}
\end{figure}

Figure~\ref{fig:calculus} shows the calculus. It differs from
\GH{Figure~5} as follows.

\begin{itemize}
\item \emph{Antecedents} $\Gamma$ are sets of judgments $Y : \phi$ about
  statement variables only; $\dom(\Gamma)$ is the set of such $Y$. Formulas
  in judgments need not be closed.
\item \emph{(Call)}: the side condition is $Y_m \notin \dom(\Gamma)$ instead
  of $Y_m : \phi_m \notin \Gamma$. With the latter, (Call) could be applied
  again to $\call{m}$ with a different contract, which the corrected
  semantics below does not support. The substitution in the premise replaces
  each call $\call{m'}$ with $Y_{m'} \in \dom(\Gamma, Y_m : \phi)$ by
  $\skipS; Y_{m'}$, as in the paper.
\item \emph{(Cons)} uses an entailment oracle $\models$, which is a parameter
  of the calculus. Soundness assumes the oracle is sound: $\phi' \models
  \phi$ implies $\fsem{\phi'} \subseteq \fsem{\phi}$.
\item \emph{(Unfold)} requires $\mu X.\phi$ to be closed, so that the plain
  substitution is semantically correct (Lemma~\ref{lem:subst}). This costs
  nothing in the paper's setting, where judgment formulas are closed.
\item \emph{Structural rules.} The paper calls the remaining logical rules
  standard and omits them. The formalization has (Ax), (Weak), (Cut),
  ($\land$R), ($\land$L) and (ConsL). Weakening and cut must preserve
  $\dom(\Gamma)$, because validity depends on $\dom(\Gamma)$ (see below).
  A left rule for $\lor$ is not sound: $I(Y) \subseteq \fsem{\phi_1 \lor
    \phi_2}$ does not imply $I(Y) \subseteq \fsem{\phi_1}$ or $I(Y) \subseteq
  \fsem{\phi_2}$.
\end{itemize}

\subsection{Validity}\label{sec:validity}

\GH{Definition~5.4} calls $\Gamma \vdash S : \phi$ valid if
$\Str_I\sem{S} \subseteq \fsem{\phi}$ for every $I$ in which all judgments
of $\Gamma$ are valid. With this definition the calculus is not sound when
$\Gamma$ is non-empty.

\begin{proposition}[Counterexample]\label{prop:counter}
  Let $T = \{m\,\{\call{m_1}\},\ m_1\,\{\skipS\}\}$ and
  $\Gamma = \{Y_{m_1} : \mathit{false}\}$. Then
  $\Gamma \vdash \call{m} : \Id\chop\Id\chop\mathit{false}$ is derivable
  (with any oracle), but it is not valid in the sense of \GH{Definition~5.4}.
\end{proposition}

\begin{proof}
  (Call) on $\call{m}$ leaves the premise
  $\Gamma, Y_m : \Id\chop\mathit{false} \vdash \skipS; Y_{m_1} :
  \Id\chop\mathit{false}$, since $Y_{m_1} \in \dom(\Gamma)$. It follows by
  (Seq), (Skip) and (Ax). But for $I(Y) = \emptyset$, the judgment in
  $\Gamma$ is valid, while $\Str_I\sem{\call{m}}$ contains the trace
  $s \cdot s \cdot s \cdot s$ and $\fsem{\Id\chop\Id\chop\mathit{false}} =
  \emptyset$.
\end{proof}

The derivation treats the call of $m_1$ as the continuation $\skipS;
Y_{m_1}$, but the semantics does not relate $Y_{m_1}$ to $m_1$. The
correction lets the antecedent determine the meaning of the procedures in
its domain.

\begin{definition}[Validity, corrected]\label{def:valid}
  $\Gamma \models S : \phi$ if for every $I$ such that
  $I(Y) \subseteq \fsem{\psi}$ for all $Y : \psi \in \Gamma$, we have
  $\Str^{\dom(\Gamma)}_I\sem{S} \subseteq \fsem{\phi}$.
\end{definition}

\begin{proposition}\label{prop:empty}
  For $\Gamma = \emptyset$, Definition~\ref{def:valid} coincides with
  \GH{Definition~5.4}.
\end{proposition}

\begin{theorem}[Soundness, \GH{Theorem~5.6}]\label{thm:sound}
  Let $T$ be variable-free and the oracle sound. Then every derivable
  sequent is valid in the sense of Definition~\ref{def:valid}.
\end{theorem}

\begin{corollary}\label{cor:sound}
  Let $T$ be variable-free and the oracle sound. If $\vdash S : \phi$ is
  derivable, then $\Str_I\sem{S} \subseteq \fsem{\phi}$ for every $I$, and
  if moreover $S$ is variable-free, then $\Ssos\sem{S} \subseteq
  \fsem{\phi}$.
\end{corollary}

The proof is by induction on derivations, with one lemma per rule. Only
(Call) is not routine. The paper's argument handles one procedure and
appeals to Beki\v{c}'s principle for several; the formalization instead does
fixed-point induction on the whole system of procedures at once.

\begin{lemma}[Fixed-point induction for (Call)]\label{lem:call}
  Let $T$ be variable-free, $m \notin D$, $m\,\{S_m\} \in T$, and
  $I' = I[Y_m \mapsto \gamma]$. If
  $\Str'\sem{S_m}_{\rho^{D \cup \{m\}}_{I'}, I} \subseteq \gamma$, then
  $\rho^D_I(m) \subseteq \dup\gamma$.
\end{lemma}

\begin{proof}[Proof sketch]
  $\rho^{D \cup \{m\}}_{I'}$ is a pre-fixed point of $H^D_I$, and its value
  at $m$ is $\dup\gamma$.
\end{proof}

For (Call), take $\gamma = \fsem{\phi}$. The premise of the rule is valid,
and in $\rho^{D \cup \{m\}}_{I'}$ the substitution $S_m[\skipS;Y_{m'}/m'()]$
does not change the semantics of $S_m$. So the hypothesis of
Lemma~\ref{lem:call} holds, and $\Str^D_I\sem{\call{m}} = \rho^D_I(m)
\subseteq \dup\fsem{\phi} = \fsem{\Id\chop\phi}$.

Finally, \GH{Corollary~5.5} holds as an equivalence.

\begin{corollary}[\GH{Corollary~5.5}]\label{cor:strongest}
  If $T$ and $S$ are variable-free, then $\models S : \phi$ if and only if
  $\fsem{\stf(S)} \subseteq \fsem{\phi}$.
\end{corollary}

\section{The derivation of Figure~6}\label{sec:example}

Let $T$ be the procedure table of \GH{Example~2.2}, with procedures
$\mathit{even}$ and $\mathit{odd}$, and let
$\phi_{\mathit{even}} = \mu X_{\mathit{even}}.\stf(\{\mathit{even}\},
S_{\mathit{even}})$ and $\phi_{\mathit{odd}} = \mu
X_{\mathit{odd}}.\stf(\{\mathit{odd}\}, S_{\mathit{odd}})$. The oracle is
semantic entailment: $\phi' \models \phi$ iff $\fsem{\phi'} \subseteq
\fsem{\phi}$.

\begin{theorem}[\GH{Example~5.2}]\label{thm:example}
  $\vdash \call{\mathit{even}} : \stf(\call{\mathit{even}})$ is derivable,
  by a derivation that follows \GH{Figure~6}. Consequently, every trace in
  $\Str_I\sem{\call{\mathit{even}}} = \Ssos\sem{\call{\mathit{even}}}$
  satisfies $\stf(\call{\mathit{even}})$.
\end{theorem}

The derivation makes the vertical dots of Figure~6 explicit. Between (Call)
and the call of $\mathit{odd}$, and again between (Call) and $Y_{even}$, it
applies (Unfold) and then (If), (Seq), (Skip) and (Assign) along the body.
Each (If) premise needs a (Cons) step: the formula has the shape
$(b \land \Id\chop\phi_1) \lor (\lnot b \land \Id\chop\phi_2)$, and
$\Id\chop\phi_1 \models \lnot b \lor \phi$ and $\Id\chop\phi_2 \models b
\lor \phi$ hold for every formula of that shape.

The two (Cons) steps shown in Figure~6 are where the paper says the formula
is ``too deep''. Both replace a strongest trace formula by its unfolding,
and Lemma~\ref{lem:unfold} shows that the two sides have the same
semantics:
\begin{itemize}
\item at $\call{\mathit{odd}}$, by part (a) with $n = \mathit{even}$ and
  Theorem~\ref{thm:stf}:
  \[
    \fsem{\stf(\call{\mathit{odd}})} = \Str_I\sem{\call{\mathit{odd}}} =
    \fsem{\stf(\{\mathit{even}\},
      \call{\mathit{odd}})[\phi_{\mathit{even}}/X_{\mathit{even}}]};
  \]
\item at $Y_{\mathit{even}}$, by part (b) with $n = \mathit{odd}$ and
  $m = \mathit{even}$:
  \[
    \fsem{\phi_{\mathit{even}}} =
    \fsem{(\mu X_{\mathit{even}}.\stf(\{\mathit{even}, \mathit{odd}\},
      S_{\mathit{even}}))[\phi_{\mathit{odd}}/X_{\mathit{odd}}]}.
  \]
\end{itemize}

\begin{remark}\label{rem:short}
  There is a shorter derivation. If (Call) on $\call{\mathit{odd}}$ uses as
  its contract the formula that unfolding $\phi_{\mathit{even}}$ produces,
  instead of $\stf(\call{\mathit{odd}})$, then neither of the two (Cons)
  steps is needed: the hypothesis for $Y_{\mathit{even}}$ then matches the
  formula for $Y_{\mathit{even}}$ directly.
\end{remark}

\begin{thebibliography}{GH}
\bibitem[GH]{GH} D.~Gurov and R.~H\"ahnle. \emph{An Expressive Trace Logic
  for Recursive Programs}.
  \url{https://arxiv.org/abs/2411.13125}.
\end{thebibliography}

\appendix
\section{The SOS as formalized}

\begin{figure}[h]
  \small
  \[
    \ottdrulesosXXskip{Skip} \qquad
    \ottdrulesosXXassign{Assign} \qquad
    \ottdrulesosXXcall{Call}
  \]
  \[
    \ottdrulesosXXseqOne{Seq$_1$} \qquad
    \ottdrulesosXXseqTwo{Seq$_2$}
  \]
  \[
    \ottdrulesosXXifOne{If$_1$} \qquad
    \ottdrulesosXXifTwo{If$_2$}
  \]
  \[
    \ottdrulesosXXtrXXend{End} \qquad
    \ottdrulesosXXtrXXstep{Step}
  \]
  \caption{The SOS (\GH{Figure~1}) and the induced traces
    (\GH{Definition~2.6}), as formalized (generated from the Ott source).}
  \label{fig:sos}
\end{figure}

\end{document}
